% BHU PYQ -- Political Science: Constitution of India-I (2025-26 NEP)
\documentclass[11pt,a4paper]{article}
\usepackage[a4paper,top=2.5cm,bottom=2.5cm,left=2.8cm,right=2.8cm,headheight=14pt]{geometry}
\usepackage{amsmath,amssymb}
\usepackage{fontspec}
\setmainfont{Kohinoor Devanagari}
\usepackage{microtype,fancyhdr,enumitem,hyperref,lastpage,parskip}

\hypersetup{colorlinks=true,urlcolor=black,linkcolor=black,pdftitle={POLMJ31: Constitution of India-I -- BHU NEP 2025-26}}

\pagestyle{fancy}\fancyhf{}
\renewcommand{\headrulewidth}{0.4pt}\renewcommand{\footrulewidth}{0.4pt}
\fancyhead[L]{\small \textit{Banaras Hindu University}}
\fancyhead[R]{\small \textit{POLMJ31: Constitution of India-I (2025-26)}}
\fancyfoot[C]{\small Page \thepage\ of \pageref{LastPage}}

\newlist{parts}{enumerate}{2}
\setlist[parts,1]{label=(\alph*),leftmargin=*,itemsep=4pt,topsep=3pt}
\setlist[parts,2]{label=(\roman*),leftmargin=*,itemsep=2pt,topsep=2pt}

\newcommand{\pts}[1]{\hfill \textit{[#1]}}

\begin{document}

\begin{center}
  {\large\bfseries BANARAS HINDU UNIVERSITY}\\[2pt]
  {\normalsize B. A. (Hons.) Semester III Examination, 2025-26 (NEP)}\\[6pt]
  \rule{0.8\linewidth}{0.5pt}\\[6pt]
  {\large\bfseries POLITICAL SCIENCE}\\[4pt]
  {\large\bfseries Paper Code: POLMJ-31 : Constitution of India-I}\\[4pt]
  {\normalsize Time: 2:30 Hours \hfill Full Marks: 70}\\[2pt]
  {\small REF NO. 2025-26/D-III/076}
\end{center}
\rule{\linewidth}{0.4pt}
\smallskip

\noindent \textit{Instructions: Attempt all questions. Questions of Sections A, B and C carry 10, 05 and 02 marks each respectively. Write your Roll No.\ at the top immediately on receipt of this question paper.}

\smallskip
\noindent \textit{सभी प्रश्नों के उत्तर दीजिए। खण्ड अ, ब तथा स के प्रश्न क्रमशः 10, 05 एवं 02 अंकों के हैं।}

\bigskip

\section*{SECTION -- A / खण्ड -- अ}
\noindent \textit{Note: Long Answer Type Questions. Answer approximately in 250 words.}\pts{3 $\times$ 10 = 30}

\noindent \textit{दीर्घ उत्तरीय प्रश्न। लगभग 250 शब्दों में उत्तर दीजिए।}

\begin{enumerate}[label=\textbf{\arabic*.},leftmargin=*,itemsep=12pt]

  \item Discuss the composition, role and significance of the Constituent Assembly of India.\pts{10}
  
  भारत की संविधान सभा की संरचना, भूमिका एवं महत्त्व की विवेचना कीजिए।

  \begin{center}\textbf{OR / अथवा}\end{center}

  Examine the philosophical foundations of the Indian Constitution as reflected in the Preamble.\pts{10}
  
  उद्देशिका (प्रस्तावना) में प्रतिबिंबित भारतीय संविधान के दार्शनिक आधारों का परीक्षण कीजिए।

  \item Critically evaluate the Fundamental Rights guaranteed under the Indian Constitution.\pts{10}
  
  भारतीय संविधान के अंतर्गत गारंटीकृत मौलिक अधिकारों का समालोचनात्मक मूल्यांकन कीजिए।

  \begin{center}\textbf{OR / अथवा}\end{center}

  Discuss the Directive Principles of State Policy and explain their relationship with Fundamental Rights.\pts{10}
  
  राज्य के नीति निर्देशक तत्त्वों की विवेचना कीजिए तथा मौलिक अधिकारों के साथ उनके सम्बन्धों की व्याख्या कीजिए।

  \item Describe the powers, functions and position of the President of India.\pts{10}
  
  भारत के राष्ट्रपति की शक्तियों, कार्यों एवं स्थिति का वर्णन कीजिए।

\end{enumerate}

\bigskip

\section*{SECTION -- B / खण्ड -- ब}
\noindent \textit{Note: Short Answer Type Questions. Answer approximately in 125 words.}\pts{4 $\times$ 5 = 20}

\noindent \textit{लघु उत्तरीय प्रश्न। लगभग 125 शब्दों में उत्तर दीजिए।}

\begin{enumerate}[label=\textbf{\arabic*.},leftmargin=*,start=4,itemsep=10pt]

  \item Write a short note on the Fundamental Duties (Article 51A).\pts{5}
  
  मौलिक कर्तव्यों (अनुच्छेद 51A) पर एक संक्षिप्त टिप्पणी लिखिए।

  \begin{center}\textbf{OR / अथवा}\end{center}

  Explain the concept of Judicial Review in India.\pts{5}
  
  भारत में न्यायिक पुनरावलोकन (Judicial Review) की अवधारणा को स्पष्ट कीजिए।

  \item Discuss the legislative powers of the Indian Parliament.\pts{5}
  
  भारतीय संसद की विधायी शक्तियों की विवेचना कीजिए।

  \begin{center}\textbf{OR / अथवा}\end{center}

  Describe the powers and position of the Prime Minister of India.\pts{5}
  
  भारत के प्रधानमंत्री की शक्तियों एवं स्थिति का वर्णन कीजिए।

  \item Write a note on the composition and jurisdiction of the Supreme Court of India.\pts{5}
  
  भारत के उच्चतम न्यायालय के गठन एवं क्षेत्राधिकार पर एक टिप्पणी लिखिए।

  \begin{center}\textbf{OR / अथवा}\end{center}

  Explain the Emergency provisions under the Indian Constitution.\pts{5}
  
  भारतीय संविधान के अंतर्गत आपातकालीन प्रावधानों की व्याख्या कीजिए।

  \item Explain the Basic Structure Doctrine propounded by the Supreme Court.\pts{5}
  
  उच्चतम न्यायालय द्वारा प्रतिपादित मूल संरचना के सिद्धान्त (Basic Structure Doctrine) की व्याख्या कीजिए।

\end{enumerate}

\bigskip

\section*{SECTION -- C / खण्ड -- स}
\noindent \textit{Note: Very Short Answer Type Questions. Answer approximately in 15--20 words.}\pts{10 $\times$ 2 = 20}

\noindent \textit{अति लघु उत्तरीय प्रश्न। लगभग 15--20 शब्दों में उत्तर दीजिए।}

\begin{enumerate}[label=\textbf{\arabic*.},leftmargin=*,start=8,itemsep=8pt]
  \item 
  \begin{parts}
    \item Which Article of the Indian Constitution guarantees the Right to Constitutional Remedies?\pts{2}
    
    भारतीय संविधान का कौन-सा अनुच्छेद संवैधानिक उपचारों के अधिकार की गारंटी देता है? (Article 32)

    \item Under which Article can National Emergency be declared?\pts{2}
    
    किस अनुच्छेद के तहत राष्ट्रीय आपातकाल की घोषणा की जा सकती है? (Article 352)

    \item Who was the Chairman of the Drafting Committee of the Constituent Assembly?\pts{2}
    
    संविधान सभा की प्रारूप समिति के अध्यक्ष कौन थे? (Dr. B. R. Ambedkar)

    \item What is the minimum age required to become the President of India?\pts{2}
    
    भारत का राष्ट्रपति बनने के लिए न्यूनतम आयु क्या है? (35 years)

    \item Name the two Houses of the Indian Parliament.\pts{2}
    
    भारतीय संसद के दोनों सदनों के नाम लिखिये। (Lok Sabha and Rajya Sabha)

    \item By which Constitutional Amendment were Fundamental Duties added?\pts{2}
    
    किस संविधान संशोधन द्वारा मौलिक कर्तव्य जोड़े गए थे? (42nd Amendment Act, 1976)

    \item What is Writ of Habeas Corpus?\pts{2}
    
    बंदी प्रत्यक्षीकरण रिट (Habeas Corpus) क्या है?

    \item Who administers the oath of office to the President of India?\pts{2}
    
    भारत के राष्ट्रपति को पद की शपथ कौन दिलाता है? (Chief Justice of India)

    \item What is the term of office of a Rajya Sabha member?\pts{2}
    
    राज्य सभा सदस्य का कार्यकाल कितना होता है? (6 years)

    \item Under which Article is the Supreme Court established as a Court of Record?\pts{2}
    
    किस अनुच्छेद के तहत उच्चतम न्यायालय को अभिलेख न्यायालय (Court of Record) बनाया गया है? (Article 129)
  \end{parts}
\end{enumerate}

\end{document}
